\section{Statistical testing protocol}
\label{app:statistical_testing}

Point estimates and significance annotations answer different questions. Forecasting and alignment entropy tables first summarize finite completed seeds within each system by \texttt{scipy.stats.trim\_mean} with proportion \(0.25\), then take an arithmetic mean across systems. This sorts \(n\) values and discards \(\lfloor n/4\rfloor\) from each tail; complete \(15\)-seed cells therefore average their central nine. We refer to this implementation as the interquartile mean (IQM). Forecasting averages all \(15\) controlled systems; primary alignment averages the \(14\) systems with at least two observed labels and retains an all-\(15\) sensitivity. The descriptive family-count column instead averages per-system seed means. We do not pool trajectories, seeds, and systems into one inferential sample.

\paragraph{Forecasting metric and missingness.}
For start \(i\), rollout step \(t\), and state coordinate \(j\), let
\[
  e_{i,t}=\sum_{j=1}^{d_x}(\hat x_{i,t,j}-x_{i,t,j})^2.
\]
For the KAE rows used in formal comparisons, the reported MSE through horizon \(H\) takes the mean of finite \(e_{i,t}\) over \(t=1,\ldots,H\) within each start, then averages finite per-start means. It is therefore a state-summed squared error, not per-coordinate MSE. A rollout can contribute through its finite prefix even if it is nonfinite before \(H\), so the KAE ranking does not enforce full-horizon finiteness. Local-linear mixture controls use the same finite-step convention; the unmatched classical DMD/EDMD evaluator requires a finite ordinary through-\(H\) mean for each contributing start, as detailed in \cref{app:experimental_details}.

The error \(E_{s,r,h}\) used below is also best-periodic: for every system \(s\), seed \(r\), model, and horizon \(h\), the lowest MSE is selected on the same test trajectories over the fixed cadence grid. This test-set oracle is shared by the point estimates and formal KAE comparisons and should not be interpreted as validation-frozen deployment selection.

\paragraph{Controlled forecasting effects.}
For candidate model \(a\), Dense MLP baseline, system \(s\), seed \(r\), and horizon \(h\), the paired effect is
\[
  \Delta_{s,r,h}^{a}
  =
  \log_{10}E^{a}_{s,r,h}
  -
  \log_{10}E^{\rm Dense}_{s,r,h}.
\]
Only mutually finite positive MSEs are eligible. Within each system and candidate--horizon cell, a one-sided paired Wilcoxon signed-rank test uses the common seeds with alternative \(\Delta<0\). At least four finite pairs are required. Raw per-system \(p\)-values are Holm-corrected across eligible systems in that candidate--horizon family. A main-table star denotes corrected \(p<0.05\); every displayed controlled sparse-row forecasting cell passes on all \(15/15\) systems.

The per-system tables in \cref{app:controlled_results} display the arithmetic mean of the paired seed effects \(\Delta_{s,r,h}^{a}\), together with the Wilcoxon/Holm \(p\)-value. The displayed effect is a mean, not the median used by an older forest-plot artifact. Because the signed-rank test depends on ranks whereas the mean does not, an extreme finite outlier can make their directions appear inconsistent; the Transfer-gated local-linear LISTA--BD H500 cell is the one such case, with \(12/15\) negative paired effects but a positive arithmetic mean.

\paragraph{Transductive alignment effects.}
The alignment partition \(F_{\rm abs}^{\rm align}\) is fit on all states in an evaluation trajectory collection and scored on states at or above the within-label \(75\)th center-margin percentile contained in that collection; ties are retained, so the subset need not contain exactly one quarter of a label, and no between-label count balancing is performed. Labels and centers are native for the two gated systems. For the \(13\) catalog systems they are nearest-center proxies obtained by deterministic clustering of long-rolled endpoints at the known benchmark count. The margin \(d_2-d_1\) is nearest-center separation, not a true boundary-distance measurement. \(H(B\mid F_{\rm abs}^{\rm align})\) uses natural logarithms and is reported in nats. The seed is the paired unit within each system and the one-sided alternative is candidate \(<\) Dense MLP. Holm correction is across eligible systems within each model--metric cell. Triple-well Duffing's frozen proxy realizes one observed label, so its zero entropy for every model is uninformative and it is excluded from primary aggregates; the all-zero paired differences and undefined signed-rank statistic are consequences. Every displayed alignment star passes on all \(14/14\) eligible systems. The observed high-margin-slice family count has no directional test because closeness to the number of represented evaluation labels, rather than monotone increase, is the relevant descriptive comparison. These tests quantify model differences in the transductive diagnostic; they do not turn it into an out-of-sample classification experiment.

\paragraph{Dysts forecasting effects.}
The historical Dysts KAE rows and their associated point estimates, sign tests, Holm corrections, and significance markers are quarantined because the KAE and standalone-control trajectories used mismatched physical cadences. No Dysts inferential result is reported until the corrected all-six-row campaign is complete.

\paragraph{Descriptive comparisons.}
The standalone DMD, EDMD, and local-linear controls configure three seeds and use independently generated evaluation trajectories; their per-system summaries omit nonfinite seed values, so they are not included in paired KAE tests. The staged support-family result reports \(225\) system--seed pairs nested within \(15\) systems, plus MSE ratios, without a confirmatory \(p\)-value. Its checkpoint score applies the same finite-prefix state-summed error reducer as the reported table, minimizes separately over the eight cadences at H100, H500, and H1000, and averages those three minima. The \(32\) selector starts are the first \(32\) of the \(100\) reported starts, whereas the global baseline selects its checkpoint on \(16\) distinct starts with an H200 final-error proxy. This \(32\%\) staged overlap and selector asymmetry make the comparison optimistic and preclude interpreting the paired win counts as held-out causal evidence. The Local-linear gates coordinate intervention uses one LISTA checkpoint at seed \(0\); bootstrap intervals summarize its \(100\) starts, while random-support medians and interquartile ranges pool \(2{,}000\) dependent point--shuffle outcomes per horizon. These describe within-checkpoint variability only.

\paragraph{High-dimensional stress tests.}
The Lorenz--96, Allen--Cahn temporal-support, and forecast-optimized
Allen--Cahn packets each average trajectories within model seed and use ten
paired model-initialization seeds as inferential units; trajectories are not
independent replicates. Lorenz--96 and the temporal-support packet use
100,000 percentile bootstrap resamples of seed-level outcomes and exact
two-sided sign-flip tests over all \(2^{10}\) assignments. Lorenz--96 applies
Holm correction over its two declared long physical horizons, \(2.5\) and
\(5\); temporal-support forecasting corrects over physical times
\(8,12,16,20\), and its final-state support gate combines frozen absolute
thresholds with a paired sparse-minus-dense ARI interval. The separate
forecast-optimized Allen--Cahn packet uses 50,000 paired-seed bootstrap
resamples of each ratio-of-arm-means reduction. Those four intervals are
marginal; the frozen primary decision is their four-cell intersection-union
gate. An exact one-sided paired sign-flip max-\(t\) calculation over the four
H160/H200 mean/terminal cells is reported only as a multiplicity sensitivity.
H160 is a prefix of H200, so these cells are dependent. Every interval
conditions on its packet's fixed generated data and does not cover new dataset
draws.

The same-checkpoint Allen--Cahn new-initial-condition check crosses three newly
generated datasets with each of the same ten paired sparse/dense checkpoint
seeds. It first averages the three datasets within checkpoint seed, then uses
the ten paired seeds for a 100,000-resample paired bootstrap interval and exact
one-sided sign-flip test over all $2^{10}$ assignments. Thus 30 seed--dataset
cells are not treated as independent replicates, and the three dataset-specific
effects are descriptive. The H200 through-horizon endpoint was frozen only
after the original four-cell gate failed, so this inference is an outcome-aware
same-checkpoint robustness check rather than independent retraining replication
or a repair of that gate.

The secondary Allen--Cahn physical-metric packet uses the identical forecasts
and the same within-checkpoint averaging over three datasets. Its sole
statistical family contains all seven H200 cumulative phase, interface, and
energy metrics. For each metric, the one-sided paired sign-flip test enumerates
all $2^{10}$ assignments of the oriented seed effects; Holm correction is over
all seven raw values. The 100,000-resample paired-seed interval is for the
oriented absolute arm-mean difference. Relative error reductions divide that
absolute difference by the observed dense arm mean; modal accuracy is instead
reported in percentage points. Instantaneous curves, terminal values, H160,
and the late half are mandatory descriptives and receive no separate
inference. Because the metrics translate an outcome-selected H200 advantage
on unchanged checkpoints, favorable results remain secondary and cannot
reclassify the original gate.

\paragraph{Holm correction and assumptions.}
For each formal family, eligible raw \(p\)-values are sorted, multiplied by the number of remaining hypotheses, made monotone by a running maximum, and clipped at one. Paired Wilcoxon tests assume the common seed differences are the relevant exchangeable units within a controlled system. All conclusions condition on the fixed benchmark roster. Evaluation-set period selection, finite-prefix omission, staged checkpoint-selection overlap, selector asymmetry, and transductive family fitting are protocol limitations independent of the multiple-comparison correction.
